\section{Threats to validity}
\label{app:threats}

We enumerate the threats to each conclusion and the design choice that addresses it, so a
reviewer can check that the negatives are not artifacts.

\paragraph{Outcome shopping.} The central threat to any negative is that the analyst kept
searching until the result appeared. We freeze every plan, model families, budgets, splits,
policies, and gates, before test dates are read, and we forbid outcome-driven tweaks. A
failure to certify an effect is recorded as such and never reported as proof of the opposite.
The soundness of a negative depends on this, and it is enforced by construction rather than by
assertion.

\paragraph{Model capacity confound.} A view could look uninformative only because the model was
too weak to use it. We hold the model family and budget constant across views and compare views
within each family, and in Q18 we add an optimizer diagnostic that varies the neural budget. The
asset-dependent direction survives the budget change, so the sufficiency negative is not a pure
underfitting artifact. Neural cells that hit their iteration cap are flagged and never used as
the sole basis for a claim.

\paragraph{Metric and policy dependence.} A decision negative could reflect one unlucky metric
or policy. Q17 uses three policies and sweeps cost and latency scenarios. The argmax negative is
absolute (even the oracle loses), the confidence policy abstains, and the scheduling policy
straddles zero across scenarios. The conclusion is stable across the policy set, not a single
strategy.

\paragraph{Aggregation artifacts.} An aggregate over heterogeneous cohorts can manufacture or
hide an effect. We report gaps under a fixed weighting against explicit baselines and decompose
each gap into mixture and within-class terms (Appendix~\ref{app:extended}). The mixture term
carries the earlier apparent depth advantage, which is why we do not report a single aggregate.

\paragraph{Data realism.} Uncalibrated receipt time, sampled-L2 depth, and quote-based utility
limit external validity: Q17 is a decision-utility statement, not realized profitability, and
Q18's small date panel rejects a simple rule rather than proving depth never helps. We state
these limits next to each claim and withhold any claim whose inputs are not certified, most
visibly the native-FIFO Q16 endpoint.

\paragraph{Generalization.} The negatives are scoped to the venues, horizons, and assets tested.
We do not claim they transfer to other markets, and the released protocol is exactly what a
reader would rerun to test transfer elsewhere.
